\section{模型假设}

基于赛题设定与本地模拟器（\texttt{jammers-py}，与官方引擎同判据）的离线复现，本文做以下
基本假设，均在实验条件内成立：

\begin{enumerate}
  \item \textbf{干扰源静态}：干扰源一经生成，位置、波束方向与接收半径在整局内不变，
        每频道至多一个源；源位于半径 $1800$ m 作业圆域内的源生成域（最远半径
        $1770$ m）。
  \item \textbf{波束几何}：定向源只在波束半角 $90^\circ$（全角 $180^\circ$，含边界）
        内有信号；接收半径 $R$ 为 $[1000,1500]$ m 内的未知常数，各源独立。
  \item \textbf{测向特性}：示向度误差 $|\Delta\theta| \le 1^\circ$，同点同频道误差固定
        （系统偏差可被多视角交会抵消）；测向耗时 $5$ s，连续测两频道切换耗时 $1$ s；
        移动速度 $5$ m/s。
  \item \textbf{清除判据}：估计点距真值 $\le 20$ m 时清除成功（耗时 $5$ s），未中耗时
        $3$ s；估计点距真值 $\le 5$ m 时可近距直接清除（不测向）。
  \item \textbf{测距能力}：机器狗对已听到频道只获得示向度（测向），不获得距离信息；
        定位完全依靠多视角示向度的交会几何。
  \item \textbf{确定性可复现}：同一随机种子下，源布局与示向度噪声完全确定，方案
        逐字节可复现（用于 A/B 对拍）。
  \item \textbf{统计性保证}：由于 $g,\theta,R$ 连续未知，检测保证为\textbf{实测统计
        性}（$422.5$ 万例蒙特卡洛 + 对抗枚举），不声称严格不漏；论文按实测听率上报。
\end{enumerate}